\documentclass[10pt]{article}
\usepackage[T1]{fontenc}
\usepackage[utf8]{inputenc}
\usepackage[french]{babel}
\usepackage{newpxtext,newpxmath}
\usepackage[a4paper,top=1.5cm,bottom=1.35cm,left=1.65cm,right=1.65cm]{geometry}
\usepackage{microtype}
\usepackage{xcolor}
\usepackage{titlesec}
\usepackage[most]{tcolorbox}
\usepackage{enumitem}
\usepackage{array}
\usepackage{booktabs}
\usepackage{fancyhdr}

\definecolor{navy}{HTML}{1F3A5F}
\definecolor{grey}{HTML}{54606E}
\definecolor{green}{HTML}{2C6E49}
\definecolor{amber}{HTML}{9C6A16}
\definecolor{red}{HTML}{A6322E}
\definecolor{plum}{HTML}{7A4EA8}
\definecolor{boxgray}{HTML}{F3F5F8}
\definecolor{boxred}{HTML}{FCF2F1}
\definecolor{boxgreen}{HTML}{F1F7F3}
\definecolor{boxplum}{HTML}{F7F3FB}
\definecolor{boxedge}{HTML}{D9E0EA}

\setlength{\parindent}{0pt}
\setlength{\parskip}{2.6pt}
\linespread{1.0}

\pagestyle{fancy}\fancyhf{}
\lhead{\small\color{grey}Scella}
\rhead{\small\color{grey}Les meilleurs couples -- 11 août 2026}
\cfoot{\small\color{grey}\thepage}
\renewcommand{\headrulewidth}{0.3pt}

\titleformat{\section}{\color{navy}\large\bfseries}{}{0pt}{}
\titlespacing{\section}{0pt}{9pt}{3pt}
\titleformat{\subsection}{\color{navy}\bfseries}{}{0pt}{}
\titlespacing{\subsection}{0pt}{7pt}{1.5pt}

\newtcolorbox{these}{colback=boxgray,colframe=boxedge,boxrule=0.5pt,arc=2.5pt,
  left=8pt,right=8pt,top=4.5pt,bottom=4.5pt}
\newtcolorbox{alerte}{colback=boxred,colframe=red!55,boxrule=0.8pt,arc=2.5pt,
  left=8pt,right=8pt,top=4.5pt,bottom=4.5pt}
\newtcolorbox{acquis}{colback=boxgreen,colframe=green!45,boxrule=0.8pt,arc=2.5pt,
  left=8pt,right=8pt,top=4.5pt,bottom=4.5pt}
\newtcolorbox{piste}{colback=boxplum,colframe=plum!45,boxrule=0.8pt,arc=2.5pt,
  left=8pt,right=8pt,top=4.5pt,bottom=4.5pt}

\newcommand{\servi}{{\color{green}\textbf{SERVI}}}
\newcommand{\latent}{{\color{plum}\textbf{LATENT}}}
\newcommand{\unfil}{{\color{amber}\textbf{À UN FIL}}}
\newcommand{\faux}{{\color{red}\textbf{FAUX}}}
\newcommand{\sz}[1]{{\small\color{grey}\textsf{#1}}}
\newcommand{\code}[1]{\texttt{\small #1}}
\newcommand{\x}{$\times$}

\begin{document}

\begin{center}
{\Huge\bfseries\color{navy} Scella}\\[4pt]
{\large\color{grey}\itshape Les meilleurs couples : où 1 + 1 fait plus que 2, et où il ne fait pas 2}\\[4pt]
{\small\color{grey} 11 août 2026 \; · \; branche \code{couche2-m0-frontier}, HEAD \code{aac2053} \; · \; 66 crates, analyse au code}
\end{center}
\vspace{2pt}

\begin{these}
\textbf{Le résultat central.} Le meilleur couple \emph{qui existe et se montre} est la \textbf{topologie
déclarée \x{} le reçu par appel, apparié injectivement} : on ne peut pas prétendre avoir consulté trois
modèles et n'en consulter qu'un. Il est servi, il mord, et il répond exactement à la question de
conformité du domaine de Kovant.

Mais le couple le plus important est ailleurs, et je l'ai d'abord manqué. La \textbf{co-signature
Ed25519 R3} -- le seul dispositif du dépôt qui donne la \textbf{non-répudiation}, donc la vérification
\emph{par un tiers sans clé} -- \textbf{existe}, et sa moitié difficile est \textbf{déjà servie} :
vérification, dégradation automatique R3$\rightarrow$R4 sur signature invérifiable, anomalies
remontées jusqu'au client. Ce qui manque est le \textbf{producteur} : rien, sur le chemin servi, ne
co-signe. \textbf{C'est le fil le plus court à haute valeur du dépôt.}
\end{these}

\section{1 · Ce qu'est un couple, et pourquoi c'est la bonne question}

Un couple n'est pas deux composants qui coexistent : c'est deux composants qui, \textbf{ensemble,
produisent une propriété vérifiable qu'aucun n'a seul}. La distinction est tranchante, et c'est tout
son intérêt -- elle élimine d'emblée les assemblages qui font bonne figure sur un schéma d'architecture
et ne prouvent rien. Un diagramme peut relier n'importe quoi à n'importe quoi ; un couple, lui, doit
survivre à la question « et qu'est-ce que la paire garantit que chaque moitié ne garantissait pas ? ».

Chaque couple est noté sur quatre axes. La \textbf{multiplicativité} : la propriété composée est-elle
réellement neuve, ou une simple juxtaposition ? Le \textbf{mordant} : existe-t-il un test qui
\emph{échoue} quand la propriété casse -- car un contrôle qui ne peut pas échouer ne prouve rien. La
\textbf{distance au servi} : câblé aujourd'hui, ou à une soudure près ? Et la \textbf{valeur dans la
conversation Kovant} : la paire répond-elle à une question qu'un acheteur conformité pose vraiment ?

Le classement qui suit contient quatre catégories, et les deux dernières ne sont pas du remplissage.
Les \textbf{couples servis} sont vérifiés dans le code de ce matin. Un couple \textbf{à un fil} a sa
moitié difficile déjà servie et lui manque une seule pièce -- c'est la catégorie la plus utile, parce
qu'elle chiffre un travail au lieu de décrire un manque. Les couples \textbf{latents} désignent des
briques réelles et testées qu'aucune ligne ne relie encore. Et les \textbf{faux couples} sont des
paires qui \emph{ressemblent} à des couples sans en être : ce sont eux qu'on survend en réunion sans le
vouloir, et les nommer vaut mieux que de les découvrir en face d'un évaluateur.

\section{2 · Tableau synoptique}

\renewcommand{\arraystretch}{1.16}
\begin{center}\small
\begin{tabular}{@{}r >{\raggedright\arraybackslash}p{5.0cm} l >{\raggedright\arraybackslash}p{7.4cm}@{}}
\toprule
& \textbf{Le couple} & \textbf{Statut} & \textbf{Ce que la paire produit, et pas les moitiés} \\
\midrule
1 & \textbf{Topologie déclarée \x{} reçu par appel} \sz{(injectif)} & \servi &
  un panel ne peut pas être \emph{prétendu} : l'appel omis casse le verdict \\
2 & \textbf{Caviardage PII \x{} scellement}, dans cet ordre & \servi &
  le sceau couvre le texte \emph{caviardé} -- donc ce qui est réellement sorti \\
3 & \textbf{Reçu v4 \x{} ancre 3a} & \servi &
  ces agents nommés, ce nombre de pas, \textbf{aucun retiré} \\
4 & Reçu \x{} coût réel par appel & \servi &
  une trace de dépense opposable, bornée en amont \\
-- & \emph{Les deux vérifieurs zéro-dépendance} & \servi &
  \emph{amplificateur} : rend les quatre précédents \textbf{croyables} \\
5 & Canari H1 \x{} participants efe-core & \servi &
  un opérateur falsifié \textbf{refuse de servir} \\
6 & \code{slm\_mediated} \x{} sa propre chute & \servi &
  la revendication embarque son falsifieur \\
\midrule
\textbf{7} & \textbf{Co-signature R3 \x{} chaîne HMAC} \sz{(vérifieur SERVI, producteur absent)} & \unfil &
  \textbf{non-répudiation} : forger exige les \textbf{deux} clés -- \emph{le fil le plus court} \\
8 & Signature de la \emph{racine} du reçu & \latent &
  l'opérateur s'engage sur la chaîne entière, contrôlable sans clé tenant \\
9 & \code{AnchorProof} \x{} Couche 3 (jalon 3b) & \latent &
  « existait au temps T », attesté hors de l'opérateur \\
\midrule
F1 & Fil B (relais SLM) \x{} souveraineté & \faux &
  egress ni gardé ni décompté ; \textbf{trois} implémentations de la règle \\
F2 & \code{config\_hash} \x{} ce qui a tourné & \faux &
  v4 scelle le \emph{déclaré}, pas le \emph{rapporté} (M4 = \code{Stub}) \\
F3 & « vrais SLM » \x{} essaim & \faux &
  fusion latente Mode-1 numérique, pas des modèles de langue \\
\bottomrule
\end{tabular}
\end{center}

\section{3 · Les couples servis}

\subsection{3.1 · Rang 1 -- Topologie déclarée \x{} reçu par appel, apparié \emph{injectivement}}

\sz{\code{efe-swarm-gateway/src/mission\_graph.rs:475} et \code{:492}}

La propriété tient en une phrase : \textbf{on ne peut pas déclarer avoir consulté trois modèles et
n'en consulter que deux}. Le test-gardien déclare \code{claude\_a}, \code{gpt\_b}, \code{gemini\_c},
n'exécute que les deux premiers, et obtient \code{fully\_proven = false} avec \code{bridge.broken = 1}
-- le déclaré-mais-absent est \code{Missing}.

Ce qui rend ce couple excellent n'est pourtant pas ce premier test, c'est le second, et il porte sur
la partie non évidente : \textbf{l'injectivité de l'appariement}. On y déclare \emph{trois appels au
même provider} et l'on ne produit qu'un seul enregistrement. Sans injectivité, les trois n\oe uds
pointeraient le même record et le reçu conclurait \code{Consistent} : un faux positif parfait, et
indétectable, puisque toutes les empreintes seraient valides. Avec l'injectivité : un n\oe ud apparie,
deux passent \code{Missing}, verdict faux. Une implémentation naïve aurait exactement ce trou -- et
c'est précisément le genre de trou qui ne se voit pas en relisant, seulement en l'attaquant.

Un détail de conception mérite d'être relevé, parce qu'il est rare : le verdict est \textbf{décomposé}.
\code{v.swarm.ok()} reste \emph{vrai} pendant que \code{fully\_proven} est \emph{faux} -- la chaîne
d'essaim est intègre, c'est le pont vers les appels externes qui est rompu. Le reçu ne dit donc pas
seulement \emph{que} quelque chose ne va pas, il dit \emph{où}. Un auditeur qui reçoit un verdict
binaire doit enquêter ; celui-ci reçoit un diagnostic.

\begin{acquis}
\textbf{Pourquoi c'est le meilleur couple servi.} Kovant construit des essaims de petits modèles pour
la conformité en achats. « Prouver que le panel a réellement eu lieu comme il a été déclaré » \emph{est}
la question de conformité de ce domaine -- pas une propriété adjacente, la question centrale. Ce couple
y répond, et il apporte son propre falsifieur.
\end{acquis}

\subsection{3.2 · Rang 2 -- Caviardage PII \x{} scellement du reçu, \textbf{dans cet ordre}}

\sz{\code{efe-swarm-gateway/src/routes/llm.rs} -- patron servi : \code{role $\rightarrow$ gate $\rightarrow$ clamp/redact $\rightarrow$ admit $\rightarrow$ scope $\rightarrow$ run $\rightarrow$ seal}}

La propriété n'est pas « on caviarde » \emph{plus} « on scelle ». C'est que \textbf{le sceau couvre le
texte caviardé}. Le code appelle \code{redact\_for\_egress} \emph{avant tout usage} de
\code{user\_text} -- avant le hash scellé, avant le calcul d'arité -- de sorte que le reçu commet ce
qui est réellement sorti. Caviarder après scellement ne vaudrait rien : on scellerait le texte en
clair et l'on jurerait ensuite l'avoir protégé. \textbf{L'ordre est le mécanisme}, et c'est pour cela
que ce couple est un vrai couple et pas une addition.

Son mordant est instructif, parce que la faille qu'il ferme était invisible. Sans la feature
\code{pii}, la ligne de caviardage \emph{disparaissait du binaire} -- sans erreur, sans avertissement,
sans trace d'exécution. Un opérateur pouvait compiler, lire « préflight frontière OK », et expédier
chaque prompt en clair à un tiers. Le correctif ne s'est pas contenté d'un \code{\#[cfg]} au site
d'appel : la porte de boot exige désormais \textbf{soit la feature, soit une reconnaissance explicite}
de l'opérateur (\code{EFE\_LLM\_ACCEPT\_CLEARTEXT\_EGRESS=true}). Le code résume l'arbitrage mieux que
je ne le ferais : le défaut \code{Off} par \emph{policy} reste légitime, c'est le \textbf{silence} qui
ne l'était pas.

La fissure, énoncée par le code lui-même et à répéter en présentation : \textbf{sans policy tenant, le
prompt sort en clair}. La frontière n'est pas souveraine, et le module le documente au lieu de le taire.

\subsection{3.3 · Rang 3 -- Le reçu v4 qui \emph{nomme} \x{} l'ancre 3a qui \emph{publie}}

\sz{\code{efe-meta-swarm-server/src/controller/iterate.rs:280} $\rightarrow$ \code{:291}}

Pris séparément, chacun a un trou que l'autre bouche exactement. Le reçu v4 nomme ses acteurs, leurs
versions déclarées et sa règle d'agrégation -- mais l'opérateur peut en publier un préfixe tronqué,
qui se re-prouve parfaitement. L'ancre garantit qu'aucun pas n'a été retiré -- mais elle ne dit rien de
qui a contribué. \textbf{Ensemble} : « ces agents nommés, dans ces versions déclarées, ont produit
exactement ce nombre de pas, et aucun n'a été retiré. »

Vérifié dans le chemin servi : l'ancre est publiée \textbf{avant} la construction de la réponse, et un
échec de publication fait échouer le RPC. L'ordre compte autant qu'au rang 2 -- une ancre publiée après
coup n'ancre rien, elle enregistre.

Son mordant est le meilleur du dépôt parce qu'il est \emph{gênant} : la démonstration montre d'abord le
bundle tronqué \textbf{passer} la re-preuve avec la clé, avant de le faire tomber sur l'ancre sans
aucune clé. Montrer sa propre garantie échouer là où elle ne s'applique pas est ce qui rend crédible
l'endroit où elle s'applique.

\subsection{3.4 · Rang 4 -- Le reçu \x{} le coût réel, appel par appel}

\sz{\code{mission\_graph.rs:369-374} ; le plafond dans \code{routes/llm.rs:249}}

Un enregistrement par appel, avec son \code{cost\_cents} scellé dans la chaîne. En amont,
\code{clamp\_budget} borne \textbf{inconditionnellement} : le cap du tenant s'il existe, \emph{puis} le
plafond serveur \emph{toujours} -- y compris pour un tenant sans politique, qui était exactement la
branche manquante. Et la conversion sature vers le budget \textbf{d'origine}, jamais vers
\code{u32::MAX} : un dépassement arithmétique ne peut pas transformer un plafonnement en autorisation.

La paire donne une trace de dépense opposable : \emph{le budget était borné en amont, et le reçu
prouve ce qui a réellement été dépensé}. Le cap seul se contourne par un bug de conversion ; le reçu
seul enregistre une dépense que rien ne bornait.

\section{4 · L'amplificateur : les deux vérifieurs à zéro dépendance}

\sz{\code{cargo tree -p verify-mini -e normal} et \code{-p verify-meta-mini} : chacun \textbf{seul}}

Ce n'est pas tout à fait un couple -- c'est ce qui rend les quatre précédents \emph{croyables}, et à ce
titre il mérite sa place. Un reçu est une affirmation. Un vérifieur auditable en un après-midi, sans
chaîne d'approvisionnement à qui faire confiance, transforme cette affirmation en fait contrôlable par
quelqu'un qui ne nous accorde rien. Réciproquement, le zéro-dépendance n'a aucune valeur détaché d'un
reçu : c'est un binaire qui ne vérifie rien. La dépendance est mutuelle, donc le couple est réel, mais
son rôle est d'\textbf{amplifier} plutôt que de produire une propriété nouvelle.

Il y en a \emph{deux}, ce qui est plus significatif qu'il n'y paraît : la contrainte a été tenue deux
fois, sur deux piles indépendantes, avec deux SHA-256 et deux HMAC écrits à la main. Ce n'est pas un
accident de dépendances, c'est une discipline.

\section{5 · Deux couples servis plus étroits}

\subsection{5.1 · Rang 5 -- Canari H1 empirique \x{} participants efe-core, fail-closed au boot}

\sz{\code{efe-meta-swarm-server/src/participant/local.rs:342}}

Un participant dont la non-expansivité est falsifiée \textbf{refuse de servir}. Deux honnêtetés portées
par le code font la valeur de ce couple plutôt qu'elles ne la diminuent. La première : c'est un
\textbf{falsifieur, jamais une preuve} -- la documentation l'écrit mot pour mot, \emph{« a pass is a
falsification failure, never a proof »}. La seconde : il a une \textbf{échappatoire déclarée},
\code{h1\_canary.mode = "off"} dans le manifeste, auquel cas le participant sert avec un avertissement
bruyant plutôt qu'en silence.

L'actif intellectuel est cependant ailleurs, et il est plus rare qu'un garde-fou : la mesure a
\textbf{réfuté le certificat analytique} du document de conception. $\hat{L} = 0{,}69$ en topologie
dynamique et $0{,}366$ en topologie fixe, contre un $\beta_1 \approx 0{,}17$ annoncé. La borne honnête
est la borne mesurée. Un dépôt qui publie la réfutation de sa propre note théorique dit quelque chose
sur la manière dont il travaille.

\subsection{5.2 · Rang 6 -- \code{slm\_mediated} \x{} sa propre chute}

\sz{\code{routes/llm.rs:166} -- \code{slm\_mediated: !hops.is\_empty() \&\& slm == hops.len()}}

La revendication « relais médié par le petit modèle » n'est vraie que si \textbf{tous} les hops sont
SLM ; une seule dégradation la fait tomber, visiblement. Et elle est dérivée de la sortie
\textbf{scellée}, jamais recalculée ailleurs -- donc la vue ne peut pas diverger de ce qui a été
attesté.

C'est le micro-patron le plus élégant du dépôt : \textbf{la revendication embarque son propre
falsifieur, et se calcule depuis la source scellée uniquement}. C'est aussi, thématiquement, le couple
le plus proche de la thèse de Kovant -- ce qui en fait un bon objet de discussion, à condition de ne
jamais le confondre avec le faux couple F1 ci-dessous.

\section{6 · Rang 7 -- la co-signature R3 : le vérifieur est servi, le producteur manque}

C'est la découverte de cette analyse, et elle m'a d'abord échappé. J'avais écrit, dans une première
version de ce classement, que le couple « reçu HMAC \x{} Ed25519 » restait à inventer. C'est faux. Il
est écrit, il est nommé \textbf{R3 / co-signature T-WITNESS}, il vit dans
\code{efe-swarm-provenance/src/cosign.rs} -- et il est conçu mieux que ce que j'avais esquissé.

\subsection{6.1 · Ce que R3 fait, et pourquoi l'asymétrie est son essence}

Un n\oe ud distant \textbf{pinné} co-signe sa propre computation avec sa \textbf{propre clé privée
Ed25519}, dont l'orchestrateur -- et tout tiers -- ne détient que la partie \textbf{publique}. Le
record passe de « \emph{un} distant a renvoyé X » à « \textbf{le n\oe ud N, dont tu as pinné la clé,
jure sous sa signature non-répudiable avoir produit X} ».

Le module explique lui-même pourquoi HMAC ne pouvait pas suffire, et c'est exactement l'argument :
une co-signature symétrique exigerait que le vérifieur \emph{détienne le secret du n\oe ud}, donc
qu'il puisse le \textbf{forger} -- pas de non-répudiation. Les deux systèmes sont \textbf{orthogonaux} :
le HMAC tenant scelle le record (co-signature \emph{incluse}) dans le \code{leaf\_hash} ; l'Ed25519 du
n\oe ud signe le tuple \code{prompt\_hash} $\Vert$ \code{output\_hash}. D'où la propriété composée, qui
est le c\oe ur du couple : \textbf{forger un R3 exige les DEUX clés}.

\subsection{6.2 · La règle de fer, et pourquoi elle fait la qualité du couple}

R3 embarque une discipline que j'ai rarement vue écrite aussi nettement, la \emph{no false elevation}
(I-RC9) : le barreau \textbf{scellé} est la \emph{revendication} du producteur, le barreau
\textbf{vérifié} est ce que la preuve cryptographique \emph{soutient}. Le vérifieur n'accepte jamais R3
sur parole -- une co-signature invérifiable (n\oe ud non pinné, révoqué, tuple altéré, signature
fausse) \textbf{redescend} le record à R4, la transaction se re-prouve toujours, et le verdict le
\textbf{signale} au lieu de se taire.

Et ce n'est pas une intention de documentation : \code{n17\_chain.rs:346} calcule
\code{(held\_strength, cosign\_anomalies)}, la ligne 365 les place dans le verdict, la ligne 393
déclenche la dégradation sur échec de \code{verify\_cosignature}. La chaîne va jusqu'au client --
\code{routes/llm.rs:201-202} sert \code{weakest\_rank} et \code{weakest\_label} dans la réponse de
mission. \textbf{La moitié difficile est servie de bout en bout.}

\begin{alerte}
\textbf{Et pourtant la propriété ne se déclenche jamais.} Aucun chemin servi ne \emph{produit} de
co-signature : les seules constructions de \code{NodeCoSignature} hors définitions de types sont des
\code{\#[test]} (\code{efe-swarm-sdk/src/n17.rs:322}), et le registre de n\oe uds pinnés n'est peuplé
que dans des tests (\code{NodeRegistry::empty().with\_node(...)}, \code{attest.rs:219}). En production
le registre est vide, aucun record n'est R3, et le barreau le plus faible est donc \textbf{R4 par
construction}. C'est un falsifieur servi qui n'a rien à falsifier.
\end{alerte}

C'est, une fois de plus, la signature du dépôt : \emph{la capacité est construite, la vérification est
câblée, le producteur manque}. Exactement la forme du trou « aucun producteur » de la pile B -- qui,
lui, a été fermé. Le précédent existe donc, et il dit que ce fil-là se pose.

\subsection{6.3 · Les deux chantiers latents qui restent, et qu'il ne faut pas confondre}

R3 donne la non-répudiation \textbf{du calcul d'un n\oe ud} : \emph{qui} a produit cette sortie. Il ne
donne pas l'engagement de l'opérateur sur \textbf{la chaîne entière}. Ce sont deux propriétés
distinctes, et les confondre serait exactement le genre de raccourci que ce document existe pour
éviter.

\textbf{Rang 8 -- signer la \emph{racine} du reçu.} Aucune route ne le fait, et les matériaux sont
pourtant là. Le crate d'identité de publieur apporte \code{ed25519-dalek} 2.1, le hachage canonique du
corps, le registre et le journal chaîné. Il est même déjà câblé ailleurs : le runner de conformité de
\code{efe-compliance-wiring} détient un vérifieur de signature réel derrière un trait
\sz{(\code{src/runner.rs:79})}, couvert par un test d'audit à contrat d'interface. Le manque est donc
sur le chemin du reçu, pas dans la brique.

\textbf{Rang 9 -- l'ancrage tiers (3b).} \code{forensic-export/src/anchor.rs:162} porte déjà
\code{verify\_anchor} et la liaison de racine décisive. L'ancre donne « existait au temps T » ; la
signature de racine donne « et cet opérateur s'y est engagé ». Ensemble, elles ferment le modèle de
confiance par ses deux bouts.

\section{7 · Les faux couples -- la section à lire deux fois}

\begin{alerte}
\textbf{F1 · Fil B (relais SLM) \x{} souveraineté. Ce n'est pas un couple, et le code l'avoue.}
\code{frontier.rs:118-133} : l'URL du reformulateur SLM \textbf{n'est pas contrainte au local}, son
egress \textbf{n'est ni gardé par l'allowlist tenant ni décompté} (le hop est budgété 0 ¢), et le
système \textbf{avertit au boot} plutôt que de prétendre une souveraineté qu'il ne vérifie pas. Pire
pour la robustesse : la règle de souveraineté existe en \textbf{trois} implémentations --
\code{provenance-capture/src/guard.rs}, \code{efe-swarm-sdk/src/backend/mod.rs}, et un « miroir
minimal » \code{slm\_url\_looks\_local} qui, délibérément, ne dépend d'aucune des deux. Elles ont déjà
divergé une fois en contournement SSRF. \textbf{« Petits modèles locaux, donc vos données restent chez
vous » serait faux aujourd'hui} -- et c'est la phrase qu'on prononce sans y penser.
\end{alerte}

\textbf{F2 · \code{config\_hash} \x{} ce qui a réellement tourné.} Le reçu v4 scelle la configuration
que le \textbf{manifeste demandait}, pas celle que le participant a \textbf{dit servir} : le certificat
porté sur le fil est encore un \code{Stub} -- un marqueur de phase 1 qui, comme son module l'écrit,
n'effectue \textbf{aucun contrôle} \sz{(\code{client/remote\_participant.rs:404})}. La
formule « cet agent, dans cette version précise » atteste donc une \emph{exigence}, pas une
\emph{réponse}. C'est la surinterprétation la plus tentante de toute l'histoire v4, parce qu'elle est
vraie à un mot près -- et que le mot compte.

\textbf{F3 · « Vrais SLM » \x{} essaim.} Les participants du mode réel exécutent une fusion latente
Mode-1 \emph{numérique} sur des templates. Appeler le méta-essaim un « essaim de SLM » serait faux, et
le dire devant Kovant serait doublement coûteux : c'est leur domaine, ils le verraient.

\section{8 · Le jugement}

\textbf{Le meilleur couple à montrer : le rang 1.} Topologie déclarée \x{} reçu apparié injectivement.
Il est servi, il mord, son falsifieur couvre le cas non évident que rate une implémentation naïve, et
il répond à la question de conformité centrale du domaine visé. Il tient dans une démonstration de dix
minutes, aujourd'hui.

\textbf{Le meilleur couple à finir : le rang 7.} La co-signature R3 est le seul dispositif du dépôt
qui produise de la \textbf{non-répudiation} -- donc la seule voie hors du caveat répété partout
ailleurs, \emph{tamper-evident pour un détenteur de clé, pas trustless pour un tiers}. Sa moitié
difficile est écrite, servie, et remonte jusqu'au client. Il lui manque un producteur.

\begin{acquis}
\textbf{La recommandation tient en une ligne : montrer le rang 1, poser le fil du rang 7.} Le premier
est ce que l'on met sur la table ; le second est le plus court chemin identifié vers un changement de
\emph{nature} -- pas d'étendue -- de ce que Scella a le droit d'affirmer. Et il a un précédent dans ce
dépôt : le trou « aucun producteur » de la pile B avait la même forme, et il a été fermé.
\end{acquis}

Un mot sur la méthode, parce qu'il porte sur ce classement lui-même. La première version de ce
document annonçait le rang 7 comme « à inventer ». Une seule commande de contrôle -- destinée à
\emph{confirmer} une absence -- a rendu trois lignes inattendues et renversé la recommandation. C'est
la raison d'être de l'annexe : les affirmations négatives sont celles qui coûtent le plus cher quand
elles sont fausses, et ce sont précisément celles qu'on ne relit pas.

\section{Annexe · Vérifier chaque couple soi-même}

Un classement de couples est une suite d'affirmations sur ce que le code fait \emph{ensemble} -- c'est
précisément le genre d'énoncé qu'on ne peut pas contrôler en lisant un fichier à la fois. Chacun est
donc ramené ici à son emplacement ou à sa commande, depuis la racine du dépôt. Les lignes \faux{} sont
les plus importantes à pouvoir rejouer : ce sont des affirmations \emph{négatives}, et une affirmation
négative ne vaut que si elle est réfutable.

\renewcommand{\arraystretch}{1.2}
\begin{center}\small
\begin{tabular}{@{}r >{\raggedright\arraybackslash}p{4.35cm} >{\raggedright\arraybackslash}p{11.0cm}@{}}
\toprule
& \textbf{Le couple} & \textbf{Où il se vérifie} \\
\midrule
1 & Topologie \x{} reçu injectif &
  \code{mission\_graph.rs:475} \sz{(appel omis)} et \code{:492} \sz{(injectivité)} \\
2 & PII \x{} scellement &
  \code{routes/llm.rs} -- \code{redact\_for\_egress} appelé \textbf{avant} le hash scellé \\
3 & Reçu v4 \x{} ancre &
  \code{controller/iterate.rs:280} \sz{(publication)} $\rightarrow$ \code{:291} \sz{(réponse)} \\
4 & Reçu \x{} coût &
  \code{mission\_graph.rs:369-374} ; le plafond dans \code{routes/llm.rs:249} \\
-- & Vérifieurs zéro-dép &
  \code{cargo tree -p verify-mini -e normal} \sz{et} \code{-p verify-meta-mini} \\
5 & Canari H1 \x{} efe-core &
  \code{participant/local.rs:342} \\
6 & \code{slm\_mediated} \x{} sa chute &
  \code{routes/llm.rs:166} \\
\textbf{7} & \textbf{Co-signature R3} \sz{(à un fil)} &
  le vérifieur servi : \code{cosign.rs:276}, \code{n17\_chain.rs:346/365/393},
  \newline client : \code{routes/llm.rs:201-202} \sz{(\code{weakest\_rank})} ;
  \newline \textbf{aucun producteur} :
  \code{grep -rn NodeCoSignature crates --include='*.rs' | grep -v test} \\
8 & Signature de racine \sz{(latent)} &
  la brique : \code{efe-compliance-wiring/src/runner.rs:79} \sz{(\code{Arc<dyn SignatureVerifier>})} \\
9 & \code{AnchorProof} \sz{(latent)} &
  \code{forensic-export/src/anchor.rs:162} $\rightarrow$ \code{pub fn verify\_anchor} \\
F1 & Fil B \x{} souveraineté &
  \code{frontier.rs:118-133} ; les trois copies :
  \code{grep -rl 'fn is\_local\_endpoint' crates} + \code{slm\_url\_looks\_local} \\
F2 & \code{config\_hash} \x{} le servi &
  \code{client/remote\_participant.rs:404} $\rightarrow$ \code{Stub} \\
-- & 515 tests, 0 échec &
  \code{cargo test -p efe-meta-verify -p verify-meta-mini -p efe-meta-swarm-server} \\
\bottomrule
\end{tabular}
\end{center}

Un avertissement sur la ligne 7, parce que c'est celle qui porte la recommandation -- et parce que
c'est celle qui m'a pris en défaut. Elle affirme deux choses de natures différentes : que le vérifieur
R3 \textbf{est servi}, contrôlable en ouvrant quatre fichiers, et qu'\textbf{aucun producteur
n'existe}, ce qui est une \emph{absence}. La seconde a été établie en écartant les fixtures de test ;
si un producteur existe ailleurs dans les 66 crates, le fil est déjà posé et la recommandation devient
sans objet. Ce serait une erreur dans le bon sens, et elle se détecte avec la commande ci-dessus.

C'est du reste ainsi que la version précédente de ce classement est tombée : la commande de la ligne 7
avait été écrite pour \emph{confirmer} que le reçu ne signait rien. Elle a rendu trois lignes, et le
rang 7 est passé de \mbox{« à inventer »} à \mbox{« à un fil »}.

\begin{these}
\small\textbf{Portée.} Ce classement a été établi contre le code du 11 août sur les Couches 2 et 3, les
deux piles de reçu et le Fil B. Il n'a \textbf{pas} examiné la couche nano ni \code{efe-gold-test-gate},
qui peuvent porter des couples non vus. La chaîne du reçu mesure \textbf{515 tests, 0 échec}
(35 binaires), par exécution réelle. Et le rappel qui ne se négocie pas : le reçu atteste
l'\textbf{intégrité} d'une suite de pas, jamais la \textbf{correction} d'une décision.
\end{these}

\end{document}
